% Revised Section 3: theoretical analysis.
% Required packages: amsmath, amssymb, amsthm, bm, mathtools.
% Suggested theorem environments:
% \newtheorem{definition}{Definition}
% \newtheorem{theorem}{Theorem}
% \newtheorem{proposition}{Proposition}
% \newtheorem{corollary}{Corollary}
% \newtheorem{remark}{Remark}

\section{From Identifiable Histories to Causally Valid Latent Dynamics}
\label{sec:theory}

The assumptions of consistency, sequential exchangeability, and
positivity concern the observed history $H_t$. A medical world model,
however, does not generally retain $H_t$ during simulation. It replaces
the history by a learned state $Z_t=E_\phi(H_t)$ and recursively applies
an action-conditioned transition. This section distinguishes two
requirements that arise from this replacement. First, the compression
must preserve the historical distinctions needed for intervention
queries. Second, the learned transition must preserve those distinctions
as the state is propagated through time. We show that neither requirement
follows from causal identification at the level of $H_t$ or from low
factual prediction error alone.

\subsection{History-level identification need not survive compression}
\label{sec:compression-exchangeability}

Suppose that the full observed history satisfies sequential
exchangeability:
\begin{equation}
    Y_{t+1:T}^{\bar a}\perp A_t\mid H_t
    \qquad\text{for every supported future schedule }\bar a.
    \label{eq:history-exchangeability}
\end{equation}
It is tempting to replace $H_t$ by an arbitrary predictive summary
$Z_t=E_\phi(H_t)$ in \eqref{eq:history-exchangeability}. Conditional
independence, however, is not generally preserved under coarsening.

\begin{theorem}[Arbitrary latent compression can destroy exchangeability]
\label{thm:compression-breaks-exchangeability}
There exists a data-generating process satisfying consistency,
positivity, and exchangeability given $H$, together with a deterministic
encoder $Z=E(H)$, such that
\begin{equation}
    Y(a)\perp A\mid H
    \qquad\text{but}\qquad
    Y(a)\not\perp A\mid Z.
    \label{eq:exchangeability-not-preserved}
\end{equation}
Consequently, exchangeability given the full history does not imply
exchangeability given an arbitrary learned representation.
\end{theorem}

\begin{proof}
Let the observed history be $H=C$, where
$C\in\{-1,+1\}$ and $\Pr(C=1)=\Pr(C=-1)=1/2$. Assign treatment according
to
\begin{equation}
    \Pr(A=1\mid C=1)=\frac12+\delta,
    \qquad
    \Pr(A=1\mid C=-1)=\frac12-\delta,
    \qquad 0<\delta<\frac12.
    \label{eq:counterexample-propensity}
\end{equation}
Both treatments have positive probability at both values of $C$, so
positivity holds. Define the potential outcome
\begin{equation}
    Y(a)=C+\beta a.
    \label{eq:counterexample-outcome}
\end{equation}
Given $H=C$, treatment is generated by an independent Bernoulli draw and
therefore $Y(a)\perp A\mid H$. Now consider the constant encoder
$E(C)=0$. Conditioning on $Z=0$ leaves $A$ associated with $C$ by
\eqref{eq:counterexample-propensity}, while $Y(a)$ remains a function of
$C$ by \eqref{eq:counterexample-outcome}. Hence
$Y(a)\not\perp A\mid Z$.
\end{proof}

Theorem~\ref{thm:compression-breaks-exchangeability} isolates a
model-induced source of causal failure. The confounder $C$ is measured in
$H$; it becomes effectively unobserved only after the encoder discards
it. This is distinct from unmeasured confounding in the original data.

\subsection{Causally sufficient latent states}
\label{sec:causal-state-definition}

We next state conditions under which the full history can safely be
replaced by a latent state. Let $\Pi_H\subseteq\mathcal A^H$ denote a
class of treatment schedules supported by the observed data.

\begin{definition}[Intervention-sufficient representation]
\label{def:intervention-sufficient}
An encoder $E:\mathcal H\rightarrow\mathcal Z$ is intervention-sufficient
for $\Pi_H$ if, whenever $E(H_t)=E(H'_t)$,
\begin{equation}
    p\!\left(Y_{t+h}^{\bar a}\mid H_t\right)
    =
    p\!\left(Y_{t+h}^{\bar a}\mid H'_t\right)
    \label{eq:intervention-sufficiency}
\end{equation}
for every $\bar a\in\Pi_H$ and every $h\leq H$.
Equivalently, for the supported query class,
\begin{equation}
    Y_{t+1:t+H}^{\bar a}\perp H_t\mid Z_t,
    \qquad Z_t=E(H_t).
    \label{eq:outcome-sufficiency}
\end{equation}
\end{definition}

One sufficient route for transferring exchangeability from $H_t$ to
$Z_t$ is that $Z_t$ is also a balancing score.

\begin{proposition}[Transfer of exchangeability to a balancing state]
\label{prop:exchangeability-transfer}
Assume \eqref{eq:history-exchangeability}, let $Z_t=E(H_t)$, and suppose
\begin{equation}
    A_t\perp H_t\mid Z_t.
    \label{eq:balancing-score}
\end{equation}
Then
\begin{equation}
    Y_{t+1:T}^{\bar a}\perp A_t\mid Z_t.
    \label{eq:latent-exchangeability}
\end{equation}
\end{proposition}

\begin{proof}
Because $Z_t$ is a deterministic function of $H_t$,
\begin{align}
p(Y^{\bar a}\mid A_t,Z_t)
&=\int p(Y^{\bar a}\mid A_t,H_t,Z_t)
          p(H_t\mid A_t,Z_t)\,dH_t\\
&=\int p(Y^{\bar a}\mid H_t)
          p(H_t\mid Z_t)\,dH_t\\
&=p(Y^{\bar a}\mid Z_t),
\end{align}
where the second line uses exchangeability given $H_t$ and the balancing
property \eqref{eq:balancing-score}.
\end{proof}

The balancing condition is sufficient but not the only way to obtain a
valid latent adjustment state. For world-model simulation, Definition
\ref{def:intervention-sufficient} is the more direct target: histories
may be merged only if they imply the same future response under every
supported treatment query.

\begin{proposition}[Irreducible effect error under latent aliasing]
\label{prop:aliasing-lower-bound}
Suppose $E(H)=E(H')=z$, but for some schedule pair
$(\bar a,\bar a')$ and horizon $h$,
\begin{equation}
    \Delta=
    \left|
    \tau_h(H;\bar a,\bar a')
    -\tau_h(H';\bar a,\bar a')
    \right|>0.
    \label{eq:aliased-effect-gap}
\end{equation}
Every effect estimator that is a function only of
$(z,\bar a,\bar a')$ then satisfies
\begin{equation}
    \max\!\left\{
    |\widehat\tau_h(z)-\tau_h(H)|,
    |\widehat\tau_h(z)-\tau_h(H')|
    \right\}
    \geq \frac{\Delta}{2}.
    \label{eq:aliasing-max-error}
\end{equation}
If $H$ and $H'$ have equal probability, their average squared effect
error is at least $\Delta^2/4$.
\end{proposition}

\begin{proof}
The estimator must output one common value $c$ for both histories. The
triangle inequality implies
$\Delta\leq|c-\tau_h(H)|+|c-\tau_h(H')|$, proving
\eqref{eq:aliasing-max-error}. For the squared-error statement, let
$u=\tau_h(H)$ and $v=\tau_h(H')$. Then
\begin{equation}
    \frac{(c-u)^2+(c-v)^2}{2}
    =\left(c-\frac{u+v}{2}\right)^2+rac{(u-v)^2}{4}
    \geq\frac{\Delta^2}{4}.
\end{equation}
\end{proof}

Thus an encoder may remain adequate for factual prediction while being
inadequate for a reusable family of intervention queries. Once
interventionally distinct histories are aliased, neither a more accurate
decoder nor a more expressive downstream transition can recover the
discarded distinction.

\subsection{Causal sufficiency must be propagated}
\label{sec:transition-closure-theory}

An intervention-sufficient state at the rollout origin is not enough.
After one model-generated transition, the simulator operates on
$\widehat Z_{t+1}=F_\theta(Z_t,A_t)$ rather than on a state freshly
encoded from the true history. We therefore require the transition to be
closed with respect to the causal representation.

\begin{definition}[Causal transition closure]
\label{def:causal-closure}
Let $Z_t=E(H_t)$. A stochastic latent transition kernel
$K^*(dz'\mid z,a)$ is closed with respect to $E$ if
\begin{equation}
    p(E(H_{t+1})\in B\mid H_t,A_t=a)
    =K^*(B\mid E(H_t),a)
    \label{eq:stochastic-closure}
\end{equation}
for every measurable $B\subseteq\mathcal Z$ and every supported
$(H_t,a)$. In a deterministic model this reduces to
\begin{equation}
    F^*(E(H_t),A_t)=E(H_{t+1}).
    \label{eq:deterministic-closure}
\end{equation}
\end{definition}

\begin{theorem}[Validity of latent g-computation]
\label{thm:latent-g-computation}
Assume consistency, positivity, and sequential exchangeability given
$H_t$. Let $E$ be intervention-sufficient for a supported schedule class
$\Pi_H$. Suppose further that:
\begin{enumerate}
    \item the controlled latent process is Markov given $(Z_t,A_t)$ and
    admits a transition kernel $K^*$ satisfying
    Definition~\ref{def:causal-closure}; and
    \item the outcome kernel is correctly represented in latent space,
    \begin{equation}
        p(Y_t\mid H_t)=p(Y_t\mid Z_t).
        \label{eq:latent-outcome-kernel}
    \end{equation}
\end{enumerate}
Then, for every $\bar a\in\Pi_H$, recursively composing $K^*$ and the
latent outcome kernel produces the same interventional outcome
distribution as the longitudinal g-formula based on the full history:
\begin{equation}
    p_{K^*}(Y_{t+1:t+H}\mid Z_t,\operatorname{do}(\bar A=\bar a))
    =p(Y_{t+1:t+H}^{\bar a}\mid H_t).
    \label{eq:latent-g-formula-equality}
\end{equation}
Consequently, contrasting any two supported latent rollouts recovers the
true longitudinal treatment-effect trajectory.
\end{theorem}

\begin{proof}
Under consistency, sequential exchangeability, and positivity, the
interventional distribution is identified by the longitudinal
g-formula. Intervention sufficiency allows the conditional future
distribution given $H_t$ to be replaced by that given $Z_t$. At the
first step, Definition~\ref{def:causal-closure} gives the same
distribution over $Z_{t+1}=E(H_{t+1})$ as the full-history transition
under $A_t=a_t$. Condition \eqref{eq:latent-outcome-kernel} gives the
same outcome distribution. Applying the same argument inductively at
$t+1,\ldots,t+H-1$ proves
\eqref{eq:latent-g-formula-equality}. Equality of treatment effects
follows by contrasting the two identified potential-outcome means.
\end{proof}

Theorem~\ref{thm:latent-g-computation} states a positive result: a
causally sufficient representation that remains closed under the
action-conditioned transition is enough to realise g-computation in a
reusable latent simulator. The theorem does not claim that a predictive
objective necessarily learns such a representation or transition.

\subsection{Approximate closure and long-horizon effect error}
\label{sec:approximate-closure}

We now consider an imperfect deterministic world model. Let
$Z^*_{t+h}(\bar a)$ denote the target causal state along a true rollout
and define
\begin{align}
    e_h(\bar a)
    &:={\left\|
        \widehat Z_{t+h}(\bar a)-Z^*_{t+h}(\bar a)
       \right\|},
    \label{eq:latent-rollout-error-revised}\\
    \epsilon_h(\bar a)
    &:={\left\|
        F_\theta(Z^*_{t+h}(\bar a),a_{t+h})
        -Z^*_{t+h+1}(\bar a)
       \right\|}.
    \label{eq:local-transition-error-revised}
\end{align}

\begin{proposition}[Propagation of approximate closure]
\label{prop:approximate-propagation}
If $F_\theta$ is $L_F$-Lipschitz in its latent argument, then
\begin{equation}
    e_h(\bar a)
    \leq L_F^h e_0
    +\sum_{k=0}^{h-1}L_F^{h-1-k}\epsilon_k(\bar a).
    \label{eq:propagation-bound-revised}
\end{equation}
\end{proposition}

\begin{proof}
Adding and subtracting
$F_\theta(Z^*_{t+k},a_{t+k})$ gives
$e_{k+1}\leq L_Fe_k+\epsilon_k$. Recursively expanding this inequality
proves \eqref{eq:propagation-bound-revised}.
\end{proof}

Let $g_\psi=q\circ G_\psi$ be the outcome decoder and define its residual
on the target causal state by
\begin{equation}
    \eta_h(\bar a)
    :=\left|
      g_\psi(Z^*_{t+h}(\bar a))-Y_{t+h}(\bar a)
      \right|.
    \label{eq:decoder-error-revised}
\end{equation}

\begin{corollary}[Effect error is controlled by both intervention branches]
\label{cor:effect-error-revised}
If $g_\psi$ is $L_G$-Lipschitz, then
\begin{equation}
\begin{split}
    |\widehat\tau_h-\tau_h|
    \leq {}&L_G\bigl[e_h(\bar a)+e_h(\bar a')\bigr]\\
    &+\eta_h(\bar a)+\eta_h(\bar a').
\end{split}
    \label{eq:effect-bound-revised}
\end{equation}
Combining \eqref{eq:effect-bound-revised} with
\eqref{eq:propagation-bound-revised} expresses long-horizon effect error
in terms of the local transition errors accumulated along both treatment
rollouts.
\end{corollary}

\begin{proof}
Apply the triangle inequality to the two potential-outcome prediction
errors and use the Lipschitz condition on $g_\psi$.
\end{proof}

This result separates two failure modes. A large $e_0$ reflects an
insufficient or inaccurately encoded initial state; large
$\epsilon_k$ reflects failure to update and propagate an otherwise useful
state. The latter is the component directly targeted by multi-horizon
latent closure training.

\subsection{Closure as a secondary selection criterion}
\label{sec:closure-selection-theory}

Equation \eqref{eq:effect-bound-revised} does not imply the unconditional
surrogate claim
\begin{equation}
    \mathcal L_{\mathrm{closure}}\downarrow
    \quad\Longrightarrow\quad
    \mathcal R_{\mathrm{effect}}\downarrow.
    \label{eq:no-unconditional-surrogate}
\end{equation}
The target encoder may itself be interventionally insufficient, the
decoder residual may increase, or the validation and target policies may
visit different regions. We therefore use closure only as a secondary
criterion among models that remain factually adequate.

Let $\theta_{\mathrm{base}}$ denote a multi-step predictive baseline and
let $H_v$ be a factual validation horizon. For a multiplicative tolerance
$\varepsilon\geq0$, define
\begin{equation}
    \Theta_\varepsilon
    :=\left\{
      \theta\in\Theta:
      \mathcal R_{\mathrm{pred},\mu}^{(H_v)}(\theta)
      \leq
      (1+\varepsilon)
      \mathcal R_{\mathrm{pred},\mu}^{(H_v)}
      (\theta_{\mathrm{base}})
    \right\}.
    \label{eq:factual-adequacy-set}
\end{equation}
We then select
\begin{equation}
    \widehat\theta
    \in\arg\min_{\theta\in\Theta_\varepsilon}
    \mathcal R_{\mathrm{cl},\mu}^{(H_v)}(\theta),
    \label{eq:closure-selection-rule}
\end{equation}
where $\mathcal R_{\mathrm{cl},\mu}^{(H_v)}$ measures the discrepancy
between recursively predicted latent states and representations of
observed future histories under the behaviour policy. Importantly,
neither \eqref{eq:factual-adequacy-set} nor
\eqref{eq:closure-selection-rule} uses counterfactual outcomes or
treatment-effect labels.

To state what this rule can guarantee in the idealised target-policy
case, define the paired closure risk
\begin{equation}
\begin{split}
    \mathcal R_{\mathrm{cl},\Pi}^{(H)}(\theta)
    :=\mathbb E_{H_t,(\bar a,\bar a')\sim\Pi}
    \left[
      \frac1H\sum_{h=1}^{H}
      \left(e_h^2(\bar a)+e_h^2(\bar a')\right)
    \right],
\end{split}
    \label{eq:paired-closure-risk}
\end{equation}
and define $\mathcal R_{\mathrm{dec},\Pi}^{(H)}$ analogously using
$\eta_h^2(\bar a)+\eta_h^2(\bar a')$.

\begin{proposition}[Closure-based upper bound within factually adequate models]
\label{prop:closure-selection-bound}
Under the assumptions of Corollary~\ref{cor:effect-error-revised},
\begin{equation}
    \mathcal R_{\mathrm{effect},\Pi}^{(H)}(\theta)
    \leq
    4L_G^2\mathcal R_{\mathrm{cl},\Pi}^{(H)}(\theta)
    +4\mathcal R_{\mathrm{dec},\Pi}^{(H)}(\theta).
    \label{eq:effect-risk-upper-bound}
\end{equation}
Consequently, among models in $\Theta_\varepsilon$ with a common upper
bound on decoder residual risk, a smaller target-policy closure risk
yields a tighter upper bound on longitudinal effect risk.
\end{proposition}

\begin{proof}
From \eqref{eq:effect-bound-revised}, write the pointwise error as the
sum of four non-negative terms. The inequality
$(x_1+x_2+x_3+x_4)^2\leq4\sum_{j=1}^{4}x_j^2$ gives
\begin{equation}
    (\widehat\tau_h-\tau_h)^2
    \leq
    4L_G^2\bigl[e_h^2(\bar a)+e_h^2(\bar a')\bigr]
    +4\bigl[\eta_h^2(\bar a)+\eta_h^2(\bar a')\bigr].
\end{equation}
Averaging over horizons, histories, and schedule pairs proves
\eqref{eq:effect-risk-upper-bound}.
\end{proof}

Proposition~\ref{prop:closure-selection-bound} justifies closure as a
diagnostic rather than as an unconditional causal surrogate. The factual
gate prevents selecting a model whose apparent latent consistency is
obtained at the cost of unacceptable observable prediction, while the
closure criterion targets a term that directly appears in an upper bound
on effect error.

\subsection{Behaviour-policy closure under treatment-policy shift}
\label{sec:policy-shift-theory}

The rule in \eqref{eq:closure-selection-rule} uses factual closure under
the behaviour policy $\mu$, whereas Proposition
\ref{prop:closure-selection-bound} involves closure under the target
schedule distribution $\Pi$. These quantities coincide only in the
absence of relevant distribution shift.

Let $\rho_\mu(z,a)$ and $\rho_\pi(z,a)$ denote finite-horizon occupancy
distributions under the behaviour and target policies. For a non-negative
local closure loss $\ell_{\mathrm{cl},\theta}(z,a)$, if
$\rho_\pi\ll\rho_\mu$ and
$w(z,a)=\rho_\pi(z,a)/\rho_\mu(z,a)$, then
\begin{equation}
\begin{split}
    \mathcal R_{\mathrm{cl},\pi}(\theta)
    &=\mathbb E_{\rho_\mu}
      [w(Z,A)\ell_{\mathrm{cl},\theta}(Z,A)]\\
    &\leq
    \sqrt{\mathbb E_{\rho_\mu}[w^2(Z,A)]}
    \sqrt{\mathbb E_{\rho_\mu}
      [\ell_{\mathrm{cl},\theta}^2(Z,A)]}.
\end{split}
    \label{eq:closure-shift-bound}
\end{equation}
Thus observed-policy closure is most informative when target schedules
are composed from well-supported local state--action relationships. If
$\rho_\pi(z,a)>0$ where $\rho_\mu(z,a)=0$, neither closure training nor
the selection rule can provide a data-driven guarantee without additional
structural assumptions.

\paragraph{Summary.}
The results establish a hierarchy. Causal identification at the
full-history level does not automatically survive latent compression
(Theorem~\ref{thm:compression-breaks-exchangeability}). A valid latent
simulator requires both intervention sufficiency and transition closure
(Theorem~\ref{thm:latent-g-computation}). Under approximate dynamics,
closure error appears explicitly in an upper bound on longitudinal
effect error (Proposition~\ref{prop:closure-selection-bound}), but
behaviour-policy closure remains only a policy-dependent proxy
\eqref{eq:closure-shift-bound}. This motivates the factual multi-horizon
latent objective introduced in Section~\ref{sec:method} and the separate
evaluation of latent capacity, rollout horizon, and held-out treatment
schedules.

